\section{Входные и выходные данные}

\subsection{Входные данные пользовательского интерфейса}

\headingtextsep

Через пользовательский веб-интерфейс в систему поступают данные, вводимые участниками рабочих пространств при выполнении повседневных и административных операций. Эти данные вводятся через формы, таблицы, панели редактирования, поисковые строки и элементы навигации.

К основным входным данным пользовательского интерфейса относятся учётные и навигационные данные, необходимые для входа в систему, выбора рабочего пространства, перехода между проектами и подтверждения прав доступа. Для локального режима разработки также допускается ввод данных, связанных с упрощённым сценарием аутентификации.

Значительную часть входных данных составляют атрибуты проектных сущностей. Для задач вводятся наименование, сроки, приоритет, статус, исполнитель, принадлежность к группе задач, родительской задаче и другим связанным объектам. Для групп задач вводятся тип группы, её наименование, статус и порядок приоритизации. Для документов и заметок вводятся заголовки, текстовое содержимое, тип записи, связи с проектными сущностями и иные служебные признаки.

Отдельной категорией входных данных являются загружаемые файловые ресурсы. Пользователь передаёт в систему имя файла, тип содержимого и бинарные данные вложения. Помимо этого, через интерфейс поступают поисковые запросы, параметры фильтрации, сортировки, ограничения по проекту и типу сущности, а также административные данные для управления участниками, агентами, учётными данными агентов и внешними интеграциями.

\textheadingsep
\subsection{Входные данные API и интеграций}

\headingtextsep

На уровне программного интерфейса в систему поступают машинно-обрабатываемые входные данные, передаваемые веб-клиентом, MCP-утилитой и внешними сервисами. Основную часть таких данных составляют структурированные HTTP-запросы с параметрами создания, чтения, изменения и поиска сущностей системы.

Для API используются идентификаторы рабочих пространств, проектов, задач, документов, заметок, файловых ресурсов, агентов и сессий. В тело запросов передаются значения полей этих сущностей, включая статусы, сроки, типы, признаки связей, параметры ограничения доступа и наборы разрешений для агентных учётных данных. При загрузке файлов дополнительно используются multipart-запросы, содержащие бинарное содержимое и метаданные файла.

Для MCP-взаимодействия входными данными выступают параметры вызова инструментов: поисковые строки, идентификаторы задач и документов, содержимое заметок, новые значения статусов и другие данные, разрешённые scope-набором агентного ключа. Для доступа к таким операциям также используются служебные данные аутентификации и ограничения области действия ключа.

От внешних интеграций в систему поступают данные о связанных артефактах и событиях, прежде всего из GitHub. Это могут быть сведения об issue, pull request, commit, branch, состоянии синхронизации, внешних идентификаторах, ссылках и метаданных, используемых для связывания внутренних сущностей системы с внешними объектами.

\textheadingsep
\subsection{Выходные данные пользовательского интерфейса}

\headingtextsep

Через пользовательский интерфейс система выводит данные, необходимые для повседневной работы участников и администраторов. Основными выходными данными являются экранные представления рабочих пространств, проектов, задач, групп задач, документов, заметок, файловых ресурсов и журналов активности.

Для задач и групп задач пользователь получает списки, карточки, иерархические представления, статусы, сроки, признаки блокировки, связи с другими сущностями и сведения об исполнителях. Для документации и заметок выводятся древовидные структуры страниц, содержимое в формате Markdown, связанные файлы, внутренние и внешние ссылки, а также история связанных изменений.

К выходным данным интерфейса также относятся результаты полнотекстового и семантического поиска, сообщения об успешных и ошибочных операциях, уведомления о состоянии действий, данные по участникам рабочего пространства, агентам, интеграциям и записям аудита. При работе с файлами система должна отображать метаданные ресурсов, превью допустимых вложений и средства получения файла на чтение или скачивание.

\textheadingsep
\subsection{Выходные данные API и сервисных процессов}

\headingtextsep

На уровне API система возвращает структурированные выходные данные в машинно-обрабатываемом виде. К ним относятся ответы с описанием созданных, изменённых или найденных сущностей, идентификаторы объектов, статусы операций, сведения о правах доступа, наборы связанных записей и данные постраничной навигации для списков.

Для запросов поиска API должно возвращать найденные задачи, документы, заметки, файловые ресурсы и иные сущности вместе с признаками релевантности, типом объекта и данными, необходимыми для перехода к найденному элементу. Для файловых операций выходными данными являются метаданные ресурса, адрес проксируемого получения файла либо временная ссылка на его безопасное скачивание.

Сервисные процессы формируют внутренние выходные данные, используемые самой системой. К ним относятся записи аудита, результаты индексации данных для поиска, сведения о сессиях работы агентов, связи между агентными сессиями и затронутыми задачами, а также данные о выполненной синхронизации с внешними интеграциями.

При работе с агентными учётными данными выходными данными являются сведения о выпущенном ключе, его статусе, области действия, сроке действия и времени последнего использования. При этом секретная часть ключа может выводиться только в момент создания и не должна сохраняться в открытом виде в последующих ответах системы.